We have compared two visual programming languages in regard to their ability to convey software design principles (modularity, instantiation, state management,
event handling) to young learners.
On a higher level, we have analyzed how two different visual programming paradigms, block-based (Scratch) and node-based
(Godot Orchestrator), emphasize different software design concepts.
The study was conducted as an educational workshop, where two student groups were taught programming principles using one of the languages each.
Data was collected using surveys, multiple-choice tests, student observation and evaluation of their created code.

Our results show that Scratch generally correlated with a better understanding of different software design principles than Godot Orchestrator.
Based on survey, multiple choice test and artifact evaluation revealed that Scratch and Orchestrator led to a good understanding of modularity.
When it came to instantiation Godot Orchestrator did not convey this well in comparison with Scratch.
Both languages seem to be viable to teach some of the selected software design princples, though.
Scratch appears easier to use than Godot Orchestrator, as it comes with a simpler user interface, more abstraction of difficult concepts and less buildt in functions.

For a beginner, Scratch might be the better choice as it abstracts low level details (e.g., vectors) that might not be needed for initial understanding.
This is especially useful when teaching more straightforward principles like modularity, however it can also be limiting when implementing more complex concepts like
finite state machines.
Godot yields more possibilities if students already have strong programming foundations, as it yields more building blocks, customizability and optimization.
Thus, it can be used to implement more difficult software patterns elegantly too.
Overall, it remains essential to introduce any language proporly, if students have not used it before.

Similar to our study, future research should consider using other visual programming languages for teaching software design.
While we relied on Godot and its community-backed plugin Orchestrator, it would be interesting to see how an industry-grade language like Unreal Engine's Blueprint~\cite{epicgames_blueprints_2025} system could support student learning.
We did not use it in our study due to limited computing resources.
Alternatively, it might be promising to compare two educational languages (e.g., Scratch, Greenfoot) that represent different visual programming paradigms.

Utting et al. \cite{utting_alice_2010} discusses and compares different visual programming languages and their suitability for teaching, but does not include any teaching activity with students.
Furthermore, we suggest comparing different visual programming paradigms within the same framework.
This could for example, using Godot to compare Orchestrator with a block-based plugin\footnote{e.g.\ \url{https://github.com/endlessm/godot-block-coding}}.
This could potentially mitigate the bias introduced by using different game engines.

Lastly, it remains essential to compare visual and textual programming languages.
Fernandes et al. \cite{fernandes_blocks_2026} conduct such a study and find that both paradigms come with their own strengths.
However, they do not consider artificial intelligence, which is now widely available and might support students with learning one or
the other paradigm better.
There might be new ways of using textual programming or visual programming languages to teach programming concepts, if every student has an individual AI tutor.
Related litererature is still spread thin, which is why we expect a lot of interesting research opportunities in that area.
As AI is increasingly used in the industry, educators should consider where and where not its use would improve student learning.
